% ------------------------------------------------------------------------
%
\begin{frame}{Inference rules}

  Let's say we want to define a function \fun{cut} that takes a
  stack~$s$ and an index~$k$ in that stack (the top has index~0) oand
  returns the prefix of the stack~$s$ of length~$k$ (if any) and the
  corresponding suffix.

  \bigskip

  For instance,
  \begin{equation*}
    \begin{array}{r@{\;}c@{\;}l@{}}
      \fun{cut}([4; 2], 0) & \rightarrow
      & (\el, [4; 2]).\\
      \fun{cut}([5; 3; 6; 0; 2], 3) & \rightarrow
      & ([5; 3; 6], [0; 2]).
    \end{array}
  \end{equation*}
  but, for the sake simplicity, nothing is said about \(\fun{cut}([0],
  7)\) and \(\fun{cut}([0], -1)\), which we write as stuck terms:
  \begin{equation*}
    \begin{array}{r@{\;}c@{}}
      \fun{cut}([0], 7) & \not\rightarrow \\
      \fun{cut}([0], -1) & \not\rightarrow
    \end{array}
  \end{equation*}

  We have:
  \begin{equation*}
    \begin{array}{@{}r@{\;}l@{\;}lr@{\;}l@{\;}l@{}}
      \fun{cut}(s, 0) & \rightarrow & (\el, s);
      & \fun{push}(x, (t, u)) & \rightarrow & (\cons{x}{t}, u).\\
      \fun{cut}(\cons{x}{s}, k+1) & \rightarrow
      & \fun{push}(x, \fun{cut}(s, k)).
    \end{array}
  \end{equation*}

\end{frame}

% ------------------------------------------------------------------------
%
\begin{frame}{Inference rules}

  When the value of a recursive call needs to be destructured, it is
convenient to use an extension of our language to avoid creating
auxiliary functions like \fun{push}:
\begin{mathpar}
\inferrule*{}{\fun{cut}(s,0) \rightarrow (\el, s)} \;\TirName{Nil}
\qquad
\inferrule
  {\fun{cut}(s, k)             \rightarrow (t, u)}
  {\fun{cut}(\cons{x}{s}, k+1) \rightarrow (\cons{x}{t}, u)}
\,\TirName{Pref}
\end{mathpar}
The new construct is called an \emph{inference rule} because it means:
`For the value of \(\fun{cut}(\cons{x}{s}, k+1)\) to be
\((\cons{x}{t}, u)\), we infer that the value of \(\fun{cut}(s,
k)\) must be \((t, u)\)'. This interpretation corresponds to an
upwards reading of the rule \TirName{Pref} (\emph{prefix}). Just as we
compose horizontally rewrite rules, we compose inference rules
vertically, stacking them, as in
\begin{mathpar}
\inferrule
  {\inferrule{
     \inferrule{\fun{cut}([0;2], 0) \rightarrow (\el, [0;2])}
               {\fun{cut}([6;0;2], 1) \rightarrow ([6], [0;2])}}
      {\fun{cut}([3;6;0;2], 2) \rightarrow ([3;6], [0;2])}}
  {\fun{cut}([5;3;6;0;2], 3) \rightarrow ([5;3;6], [0;2])}
\end{mathpar}

\end{frame}
